\documentclass[11pt,a4paper]{article}
\usepackage[utf8]{inputenc}
\usepackage[ngerman]{babel}
\usepackage[T1]{fontenc}
\usepackage{amsmath}
\usepackage{amsfonts}
\usepackage{amssymb}
\usepackage{graphicx}
\usepackage{listings}
\usepackage{xcolor}
\usepackage{hyperref}

\title{Detaillierte RTOS-Dokumentation: Von HardFaults zur Stabilität}
\author{Danielou Mounsande}
\date{28. April 2026}

\definecolor{mGreen}{rgb}{0,0.6,0}
\definecolor{mGray}{rgb}{0.5,0.5,0.5}
\definecolor{mPurple}{rgb}{0.58,0,0.82}
\lstset{
    commentstyle=\color{mGreen},
    keywordstyle=\color{magenta},
    numberstyle=\tiny\color{mGray},
    stringstyle=\color{mPurple},
    basicstyle=\footnotesize\ttfamily,
    breaklines=true,
    captionpos=b,
    keepspaces=true,
    numbers=left,
    showspaces=false,
    showstringspaces=false,
    showtabs=false,
    tabsize=2,
    language=C
}

\begin{document}

\maketitle
\tableofcontents
\newpage

\section{Einleitung}
Diese Dokumentation dient als persönliches Logbuch und technisches Referenzwerk für die Entwicklung eines eigenen Echtzeitbetriebssystems auf dem STM32L475. Sie dokumentiert nicht nur den Code, sondern auch den Lernprozess durch gezielte Fragestellungen.

\section{Hardware-Konfiguration}
\begin{itemize}
    \item \textbf{MCU:} STM32L475VGT6 (Cortex-M4).
    \item \textbf{LEDs:} LD1 (PA5), LD2 (PB14), LD3 (PC9).
    \item \textbf{Debugger:} J-Link (GDB Server).
\end{itemize}

\section{Persönliches FAQ: Fragen und Antworten zur Architektur}

\subsection{Was war noch mal das Link Register (LR) und der Stack Pointer (SP)?}
\textbf{Antwort:} 
\begin{itemize}
    \item \textbf{SP (R13):} Zeigt auf die aktuelle Adresse im Stapelspeicher (Stack). Im Cortex-M4 verwenden wir hier den MSP (Main Stack Pointer).
    \item \textbf{LR (R14):} Speichert die Rücksprungadresse nach einem Funktionsaufruf (\texttt{BL}). In einer Exception enthält LR den magischen Code \texttt{EXC_RETURN} (z.B. \texttt{0xFFFFFFF9}).
\end{itemize}

\subsection{Warum nutzen wir nur den MSP und nicht PSP?}
\textbf{Antwort:} Gemäß den Projektanforderungen und der Vereinfachung nutzen wir nur den MSP. In komplexeren Systemen trennt man MSP (Kernel) und PSP (User Tasks), um den Kernel vor Stack-Overflows der Tasks zu schützen.

\subsection{Was passiert beim Kontextwechsel hardwareseitig?}
\textbf{Antwort:} Wenn ein PendSV-Interrupt ausgelöst wird, sichert die Hardware automatisch \texttt{xPSR, PC, LR, R12, R3, R2, R1, R0}. Wir müssen nur noch \texttt{R4-R11} manuell sichern.

\section{Forensische Analyse: Die großen Probleme}

\subsection{Der INVPC (0x00040000) Absturz}
\textbf{Problem:} Der Prozessor ging in einen HardFault mit dem Fehler "Invalid PC load".
\textbf{Ursache:} 
1. Der \texttt{BL}-Befehl im Handler überschrieb das LR (EXC_RETURN).
2. Das LSB (Least Significant Bit) des Task-PC war 0 (ARM-Modus statt Thumb-Modus).
\textbf{Lösung:} Sichern des LR in \texttt{R4} (callee-saved) und Erzwingen von \texttt{address | 0x01}.

\subsection{Das Ansaugen der Vector Table (Null-Pointer)}
\textbf{Problem:} GDB zeigte, dass \texttt{R4-R7} plötzlich Werte wie den Reset-Handler enthielten.
\textbf{Ursache:} Der Scheduler wurde durch den SysTick zu früh aufgerufen (während \texttt{HAL_Init}). Da noch keine Tasks existierten, gab der Scheduler \texttt{0x00000000} zurück.
\textbf{Lösung:} Einführung der \texttt{Global_OsIsRunning} Strategie (Boot Lock).

\section{Die endgültige Strategie: Hybride C/Assembler Architektur}
Um die Instabilität von \texttt{naked}-Funktionen zu umgehen, delegieren wir die Pointer-Mathematik an C:
\begin{lstlisting}
// C-Teil: Berechnet den SP und gibt ihn in R0 zurück
uint32_t Kernel_ContextSwitch(uint32_t current_sp) {
    if (Global_OsIsRunning) {
        Global_PointerToCurrentlyRunningTCB->u32TaskSP = current_sp;
        // ... Task Auswahl Logik ...
    }
    return Global_PointerToCurrentlyRunningTCB->u32TaskSP;
}
\end{lstlisting}
\begin{lstlisting}
// Assembler-Teil: Nur Register-Sicherung
stmdb r0!, {r4-r11}   // Sichern
bl Kernel_ContextSwitch // Berechnen in C
ldmia r0!, {r4-r11}   // Wiederherstellen
\end{lstlisting}

\section{Aktueller Status}
\textbf{Meilenstein erreicht:} Der Kernel läuft präemptiv. LD1 und LD2 blinken stabil. Der "Boot-Lock" schützt die Initialisierungsphase.

\section{Referenzen}
\begin{itemize}
    \item PM0214: STM32 Cortex-M4 Programming Manual.
    \item AAPCS: Procedure Call Standard for the Arm Architecture.
    \item Qing Li: Real-Time Concepts for Embedded Systems.
\end{itemize}

\end{document}
